\section*{第 \ref{ch:g12:contdist} 章——连续型随机变量}

\begin{solution}{exo:g12:contdist:1}
$\P(X \leq 5) = \frac{5}{15} = \frac13$;
$\P(4 \leq X \leq 10) = \frac{6}{15} = \frac25$;
$\E(X) = \frac{15}{2} = 7.5$ 分钟；
$\sigma(X) = \frac{15}{\sqrt{12}} = \frac{5\sqrt3}{2} \approx 4.3$ 分钟。
\end{solution}

\begin{solution}{exo:g12:contdist:2}
$f \geq 0$ 且
$\int_0^2 \frac38 t^2\,\dd t = \frac38\left[\frac{t^3}{3}\right]_0^2
= \frac38 \times \frac83 = 1$：是一个\omterm{def:g12:contdist:density}{概率密度}。
\[
\E(X) = \int_0^2 \frac38 t^3\,\dd t = \frac38 \times \frac{16}{4} = \frac32,
\qquad
\P(X \geq 1) = \int_1^2 \frac38 t^2\,\dd t = \frac{8 - 1}{8} = \frac78 .
\]
\end{solution}

\begin{solution}{exo:g12:contdist:3}
\emph{1.} $\E(X) = \frac{1}{0.2} = 5$ 年；
$\P(X > 5) = \eu^{-0.2\times5} = \eu^{-1} \approx 0.37$.

\emph{2.} 由无记忆性（\cref{thm:g12:contdist:memoryless}）可得，
$\pcond{X > 3}{X > 8} = \P(X > 5) = \eu^{-1} \approx 0.37$：已服役的三年
并不改变任何结果。
\end{solution}

\begin{solution}{exo:g12:contdist:4}
$\P(X > m) = \eu^{-\lambda m} = \frac12$ 给出
$m = \frac{\ln 2}{\lambda} \approx \frac{0.69}{\lambda}$，小于
$\E(X) = \frac1\lambda$。\omterm{def:g12:contdist:expo}{指数分布}的\omterm{def:g12:contdist:density}{概率密度}具有较长的右尾：少数异常长的寿命
把\emph{\omterm{def:g11:stat:mean}{平均数}}拉到
\emph{\omterm{def:g11:stat:median}{中位数}}之上，即总体中有一半超过的那个值。（这正是
\cref{ex:g12:exp:halflife}中的半衰期：一半原子能
存活过该时间，尽管平均寿命更长。）
\end{solution}

\begin{solution}{exo:g12:contdist:5}
$168 = \mu - \sigma$ 且 $182 = \mu + \sigma$：约 $68\%$。
$189 = \mu + 2\sigma$：其上方占剩余部分
$100 - 95.4 = 4.6\%$ 的一半，因此约为 $2.3\%$。
$154 = \mu - 3\sigma$：约为 $\frac{100 - 99.7}{2} = 0.15\%$。
\end{solution}

\begin{solution}{exo:g12:contdist:6}
$\P(X < 500) \leq 2.5\%$ 意味着 $500$ 至少要比均值低
两个\omterm{def:g11:stat:variance}{标准差}（正态\omterm{def:g12:randvar:rv}{分布}的 $2.3\% \approx 2.5\%$ 落在
$\mu - 2\sigma$ 以下；若取常用值 $1.96$，推理完全相同）：
\[
\mu - 2\sigma \geq 500 \iff \mu \geq 500 + 8 = 508 \text{ 克}.
\]
机器必须平均设定为 $508$ 克——这一保证的
代价是每袋多出 $8$ 克的“免费”产品。
\end{solution}

\begin{solution}{exo:g12:contdist:7}
\emph{1.} 取 $p = \frac16$，$n = 1200$：
$\sqrt{\frac{p(1-p)}{n}} = \sqrt{\frac{5/36}{1200}} \approx 0.0108$，因此
$95\%$ 的波动\omterm{def:g10:numbers:interval}{区间}为
\[
\frac16 \pm 1.96 \times 0.0108 \approx \intcc{0.146}{0.188}.
\]

\emph{2.} 观测\omterm{def:g10:stats:series}{频率} $\frac{260}{1200} \approx 0.217$ 远远落在
\omterm{def:g10:numbers:interval}{区间}之外：在 $5\%$ 水平下拒绝公平性假设。
偏差约为 $4.6$ 个\omterm{def:g11:stat:variance}{标准差}——若用正态近似，
\omterm{def:g10:proba:distribution}{概率}的量级为 $10^{-6}$，远胜于
比内–切比雪夫不等式给出的 $\leq 4.6\%$ 的上界
（见\cref{exo:g12:sums:7}）。
\end{solution}

\begin{solution}{exo:g12:contdist:8}
\emph{1.} 误差范围 $\frac{1}{\sqrt n} \leq 0.02$ 需要
$n \geq 2500$ 人。

\emph{2.} 要区分相差 $1$ 分的两个得分，误差范围必须
远小于 $0.5$ 分，按简化公式需要 $n \geq \left(\frac{1}{0.005}\right)^2 =
40\,000$ 人，对大多数民意调查来说已不切实际——
更糟的是，统计误差范围只反映\emph{抽样}误差；
系统性偏差（\omterm{def:g10:proba:sample}{样本}缺乏代表性、无应答、临近投票时的
摇摆等）不会随 $n$ 增大而缩小，通常超过 1 分。
相差 $1$ 分的竞争确实难分胜负。
\end{solution}

\begin{solution}{exo:g12:contdist:9}
\emph{1.} 由于 $\alpha > 0$，$c \leq \alpha X + \beta \leq d
\iff \frac{c - \beta}{\alpha} \leq X \leq \frac{d - \beta}{\alpha}$，因此
\[
\P(c \leq Y \leq d)
= \int_{(c-\beta)/\alpha}^{(d-\beta)/\alpha} f(t)\,\dd t .
\]
作代换 $y = \alpha t + \beta$（即在 $y$
变量下读取面积，$t = \frac{y - \beta}{\alpha}$，$\dd t = \frac{\dd y}{\alpha}$）
于是化为
$\int_c^d \frac{1}{\alpha} f\left(\frac{y-\beta}{\alpha}\right)\dd y$：
\omterm{def:g11:func:function}{函数} $g(y) = \frac1\alpha f\left(\frac{y-\beta}{\alpha}\right)$ 是
$Y$ 的一个\omterm{def:g12:contdist:density}{概率密度}。

\emph{2.} 取 $f = \varphi$，$\alpha = \sigma$，$\beta = \mu$：
\[
g(y) = \frac{1}{\sigma}\,\varphi\!\left(\frac{y - \mu}{\sigma}\right)
= \frac{1}{\sigma\sqrt{2\pi}}\,
\exp\left(-\frac{(y-\mu)^2}{2\sigma^2}\right),
\]
因此它恰是
$\mathcal N(\mu, \sigma^2) = \mu + \sigma\,\mathcal N(0,1)$ 的\omterm{def:g12:contdist:density}{概率密度}。
\end{solution}

\begin{solution}{pb:g12:contdist:1}
\textbf{1.} $\P(\text{等待} \leq 10) = \frac{10}{30} =
\frac13$；$\E = 15$ 分钟；
$\sigma = \frac{30}{\sqrt{12}} \approx 8.7$ 分钟。

\textbf{2.} $\int_0^1 2x\,\dd x = 1$：是一个\omterm{def:g12:contdist:density}{概率密度}。
$\E(X) = \int_0^1 2x^2\,\dd x = \frac23$；
$\P\left(X > \frac12\right) = \int_{1/2}^1 2x\,\dd x =
1 - \frac14 = \frac34$。

\textbf{3.} $\P(T > 15) = \eu^{-15/10} \approx 0.22$。\omterm{def:g11:stat:median}{中位数}由
$\eu^{-m/10} = \frac12$ 解得 $m = 10\ln 2 \approx 6.9$ 分钟——
小于\omterm{def:g11:stat:mean}{平均数} $10$，这是因为\omterm{def:g12:contdist:expo}{指数分布}右侧的长尾
（极少出现的超长等待）把\omterm{def:g11:stat:mean}{平均数}拉高，而使
\omterm{def:g11:stat:median}{中位数}仍停留在“典型”等待时间上。

\textbf{4.} 由无记忆性，
$\P(T > 35 \mid T > 20) = \P(T > 15) \approx 0.22$：
已等待的 20 分钟毫无作用——
这股到达流不会“变老”。

\textbf{5.} (a) \omterm{def:g12:contdist:expo}{指数分布}（放射性衰变正是为
无记忆性而设的模型）；(b) 在
$\intcc{0}{8}$ 上的\omterm{def:g12:contdist:uniform}{均匀分布}（带随机偏移的时刻表）；(c)
\omterm{def:g12:contdist:expo}{指数分布}；(d) 两者都不是：磨损的灯泡在下一小时内\emph{更}
容易失效（相比新灯泡），因此
$\P(T > s + t \mid T > s) < \P(T > t)$：老化与
无记忆性定理相矛盾。

\textbf{6.} $68\,\%$、$95\,\%$、$99.7\,\%$：这三个
标志性\omterm{def:g10:numbers:interval}{区间}。

\textbf{7.} $z = \frac{190 - 172}{8} = 2.25$：上侧
尾部\omterm{def:g10:proba:distribution}{概率} $\approx 1.2\,\%$。

\textbf{8.} $z = 2$：约 $2.3\,\%$——大约每
$44$ 人中有一人。

\textbf{9.} 规格设在 $\pm 2\sigma$ 处：约
$95.4\,\%$ 合格，$4.6\,\%$ 被拒——大规模生产时代价高昂。在
$\pm 6\sigma$ 时，失效率降至
约\emph{十亿}分之二：六西格玛换来了
批量生产而无需逐件检验的资格。

\textbf{10.} $\sigma = 5$；取\omterm{def:g12:limcont:continuity}{连续}性修正，
$\P \approx \P\left(\abs Z \leq \frac{5.5}{5}\right) =
\P(\abs Z \leq 1.1) \approx 0.73$。该定理平滑了
\cref{pb:g11:binom:1}中的高尔顿板：阶梯状的槽位
变成了\omterm{def:g12:limcont:continuity}{连续}的钟形曲线。

\textbf{11.} $n \approx \frac{1.96^2 \times 0.25}{0.005^2}
\approx 38\,400$ 人——而在这样的规模下，
\emph{非}抽样误差（抽样框、拒访、谎报）远超
误差范围：一场仅差一个百分点的竞选已超出诚实民调的能力，这也是
为什么严肃的机构此时会说“难分胜负”。

\textbf{12.} 无记忆性：自你到达之时起，剩余的
等待时间服从\omterm{def:g11:stat:mean}{平均数}为 $10$ 的\omterm{def:g12:contdist:expo}{指数分布}——即完整的
$10$ 分钟，而非 $5$ 分钟。时刻表中的“平均间隔 $10$”
容易引起误解：你的等待时间与一整段间隔服从同一\omterm{def:g12:randvar:rv}{分布}。

\textbf{13.} 随机到达的乘客落入某段间隔的\omterm{def:g10:proba:distribution}{概率}与该间隔的
长度成正比：$\P(\text{长间隔}) = \frac{15}{20} = \frac34$。
\omterm{def:g12:randvar:exp}{期望}等待时间：$\frac34 \times \frac{15}{2} + \frac14 \times
\frac52 = 5.625 + 0.625 = 6.25$ 分钟——比朴素估计的
$5$ 分钟多，尽管平均间隔为 $10$。罪魁祸首是\emph{长度偏倚
抽样}——长间隔所捕获的乘客更多。

\textbf{14.} 向后与向前的等待时间各自服从
\omterm{def:g11:stat:mean}{平均数}为 $10$ 的\omterm{def:g12:contdist:expo}{指数分布}（无记忆性从你到达的
那一刻起在两个方向上同时成立）：包含你的那段间隔平均为 $20$ 分钟
——是典型间隔的两倍。考察悖论一言以蔽之：
\emph{你恰好考察的\omterm{def:g10:numbers:interval}{区间}并非典型\omterm{def:g10:numbers:interval}{区间}，
因为你更可能落入一个大\omterm{def:g10:numbers:interval}{区间}，
即落在长间隔之中。}

\textbf{15.} 班级平均人数：$\frac{9 \times 10 + 110}{10} =
20$ 人。学生眼中的平均班级人数：
$\frac{90 \times 10 + 110 \times 110}{200} = \frac{13\,000}
{200} = 65$ 人。宣传册印的是 $20$；五分之三的学生
坐在 $110$ 人的班级中，体验的却是 $65$ 人的规模。

\textbf{16.} 受欢迎的人出现在许多人的好友名单上，因此
随机抽取“一位朋友”存在长度偏倚，偏向善于交际的人
——你的朋友正是从社交
时刻表的长间隔中抽取而来。

\textbf{17.} 对 $t \geq 0$：
$\P(T > t) = \P(-10\ln U > t) = \P\left(U <
\eu^{-t/10}\right) = \eu^{-t/10}$：恰好是\omterm{def:g12:contdist:expo}{指数分布}的
生存\omterm{def:g11:func:function}{函数}——一次取对数，便把计算机给出的
均匀噪声化为任意的等待时间。

\textbf{18.} 每个均匀随机数的\omterm{def:g11:stat:mean}{平均数}为 $\frac12$，\omterm{def:g12:randvar:exp}{方差}为
$\frac{1}{12}$：十二个之和的\omterm{def:g11:stat:mean}{平均数}为 $6$，\omterm{def:g12:randvar:exp}{方差}为
$1$——因此该算法输出的\omterm{def:g11:stat:mean}{平均数}为 $0$，\omterm{def:g12:randvar:exp}{方差}为 $1$。十二个
\omterm{def:g12:sums:indep}{独立}小扰动之和：这就是高尔顿机制，借
棣莫弗–拉普拉斯定理的庇佑，直方图已经
在肉眼看来呈钟形。

\textbf{19.} 关于模型：市场并非许多
\emph{\omterm{def:g12:sums:indep}{独立}}的小因素之和——恐慌使
一切彼此相关（上一题中 2008 年的教训），
并产生了任何正态密度都无法刻画的厚尾。当数据
仅轻声暗示 $10^{-137}$ 这样的\omterm{def:g10:proba:distribution}{概率}时，统计学家的信条是：
\emph{被拒绝的是模型，而不是这一天}。

\textbf{20.} \omterm{def:g12:contdist:uniform}{均匀分布}：各位置等可能，即诚实的“我
只知道边界”。\omterm{def:g12:contdist:expo}{指数分布}：永不
老化的危险率——衰变、到达、点击。正态\omterm{def:g12:randvar:rv}{分布}：民主的
钟形，源于众多微小\omterm{def:g12:sums:indep}{独立}因素——身高、误差、
\omterm{def:g11:stat:mean}{平均数}。陷阱在于：你抽到的东西会因你
碰到它的方式而产生偏倚——公交车、班级、朋友。全书的主线：
六年级时数果汁盒的那个孩子，如今经过十册习题，
已用一条曲线度量人群；$\varepsilon$、\omterm{def:g12:integ:area}{积分}与
中心极限定理正等待着大学分册，
来解释\emph{为什么}钟声为
万物而鸣。
\end{solution}
